%%%%%%%%%%%%%%%%%%%%%%%%%%%%%%%%%%%%%%%%%%%%%%%%%%%%%%%%%%%%%%%%%%%%%
%%%%%%%%%%%%%%%%%%%%%%%%%%%%%%%%%%%%%%%%%%%%%%%%%%%%%%%%%%%%%%%%%%%%%
%%%%%%%%%%%%%%%%%%%%%%%%%%%%%%%%%%%%%%%%%%%%%%%%%%%%%%%%%%%%%%%%%%%%%
%%%%%%%%%%%%%%%%%%%%%%%%%%%%%%%%%%%%%%%%%%%%%%%%%%%%%%%%%%%%%%%%%%%%%
\chapter{Codecs}
%%%%%%%%%%%%%%%%%%%%%%%%%%%%%%%%%%%%%%%%%%%%%%%%%%%%%%%%%%%%%%%%%%%%%
%%%%%%%%%%%%%%%%%%%%%%%%%%%%%%%%%%%%%%%%%%%%%%%%%%%%%%%%%%%%%%%%%%%%%
%%%%%%%%%%%%%%%%%%%%%%%%%%%%%%%%%%%%%%%%%%%%%%%%%%%%%%%%%%%%%%%%%%%%%

A codec\index{codec} is a software component that provides two essential functions for working with a signal: encoding and decoding. The \textbf{encoder} converts an analogue or digital image from its original (raw) format into another format with suitable parameters (encoded format). In computer networks, the size of the transmitted data is an important parameter. Reducing the size comes at the cost of lower quality and increased latency or jitter during encoding. The \textbf{decoder} converts encoded data back to its original form. In videoconferencing, encoding and decoding apply to signals of an analogue nature, namely audio and video. We therefore distinguish between video and audio codecs.

%%%%%%%%%%%%%%%%%%%%%%%%%%%%%%%%%%%%%%%%%%%%%%%%%%%%%%%%%%%%%%%%%%%%%
%%%%%%%%%%%%%%%%%%%%%%%%%%%%%%%%%%%%%%%%%%%%%%%%%%%%%%%%%%%%%%%%%%%%%
\section{Why codecs are needed}
%%%%%%%%%%%%%%%%%%%%%%%%%%%%%%%%%%%%%%%%%%%%%%%%%%%%%%%%%%%%%%%%%%%%%
%%%%%%%%%%%%%%%%%%%%%%%%%%%%%%%%%%%%%%%%%%%%%%%%%%%%%%%%%%%%%%%%%%%%%

We will explain why image encoding is necessary and why it is not sufficient to use the signal in its original form. A similar example can also be given for audio transmission.

In the real world, an image is captured several times per second and divided into frames. We denote this rate by $fps$ (frames per second). Each frame consists of pixels, where each pixel is represented as a tuple of bits. This combination represents a particular colour in the RGB colour space.

\begin{quote}
Example: For simplicity, assume that we have an image with CIF resolution of 352x288 pixels, a 24-bit colour depth and a frame rate of 20 fps. (24-bit RGB provides $256=2^8$ possible values for red, $256=2^8$ values for green and $256=2^8$ values for blue, giving $2^{24}=16777216$ different colours.) The resulting data rate for one such signal would be:

width $\times$ height $\times$ depth $\times$ fps =
$352\times288\times24\times20$ =
48660480b/s =
\textbf{5.7MiB/s}
\end{quote}

This value is far beyond the capabilities of current networks, especially when we consider that several video streams must be transmitted simultaneously during a conference. For D1 PAL resolution, the value is as high as \textbf{23MiB/s}.

%%%%%%%%%%%%%%%%%%%%%%%%%%%%%%%%%%%%%%%%%%%%%%%%%%%%%%%%%%%%%%%%%%%%%
%%%%%%%%%%%%%%%%%%%%%%%%%%%%%%%%%%%%%%%%%%%%%%%%%%%%%%%%%%%%%%%%%%%%%
\section{Video and audio transmission using the RTP protocol}
%%%%%%%%%%%%%%%%%%%%%%%%%%%%%%%%%%%%%%%%%%%%%%%%%%%%%%%%%%%%%%%%%%%%%
%%%%%%%%%%%%%%%%%%%%%%%%%%%%%%%%%%%%%%%%%%%%%%%%%%%%%%%%%%%%%%%%%%%%%

This section presents selected codecs used by contemporary videoconferencing tools. Most of these systems use the RTP/RTCP protocol for transmission, as specified in RFC 3551 \cite{RFC-3551-RTP-Profile-AudioVideo}. The video codecs include CelB, JPEG (Motion JPEG), H.261, H.263, H.263-1998 (H.263+), MPV, MP2T and nv. The audio codecs include DVI4, G722, G723, G726-40, G726-32, G726-24, G726-16, G728, G729, G729D and G729E, GSM, GSM-EFR, L8, L16, LPC, MPA, PCMA, PCMU, QCELP, RED and VDVI.

Another important document is RFC 2198 \cite{RFC-2198-RTP-Payload-Redundant-Audio}, which describes a method for transmitting audio packets over networks where packet loss may occur. To prevent interruptions in the audio, several redundant encodings are used in addition to the primary encoding.

\subsection{Video: NV}\label{kap:NV}
A historical video codec from Xerox PARC, implemented, for example, in the VIC application developed at UCL.

\subsection{Video: H.261}\label{kap:H.261}
In 1990, the H.261 low-latency codec was standardised. It is still in use today and was designed for transmitting images at QCIF and CIF resolutions (Figure \ref{fig:Resolutions}). Figure \ref{fig:Comparison-H261-NV} shows its improved quality-to-bitrate ratio compared with the older NV codec. H.261 is supported by a wide range of applications (VIC, MPlayer, VLC and FFmpeg), mainly because of its compatibility with older systems. The transmission of H.261-encoded video is also covered by RFC 2032 \cite{RFC-2032-RTP-Payload-Format-H.261}.

\begin{figure}[htp]
\centering
\includesvg[width=0.5\textwidth]{h261-vs-nv}
\caption{Comparison of the H.261 and NV codecs (source: \protect\cite{Comparison-H261-NV}).}
\label{fig:Comparison-H261-NV}
\end{figure}

\subsection{Video: H.263}\label{kap:H.263}
The H.263 codec is the natural successor to H.261. It became a standard in 1996. It combines features of H.261 with MPEG technology and optimises them for low bitrates. Compared with its predecessor, it is much more efficient at all bitrates, as shown in Figure \ref{fig:Comparison-H263-H261}. It also supports more image sizes: Sub-QCIF, QCIF, CIF, 4CIF and 16CIF. The transmission of H.263-encoded video is covered by RFC 2190 \cite{RFC-2190-RTP-Payload-Format-H.263}.

\begin{figure}[htp]
\centering
\includesvg[width=0.5\textwidth]{h263-vs-h261}
\caption{Comparison of the H.261 and H.263 codecs (source: \protect\cite{Comparison-H263-H261}).}
\label{fig:Comparison-H263-H261}
\end{figure}

\subsection{Video: H.264 (MPEG-4 AVC)}\label{kap:H.264}
One of the newer codecs intended for videoconferencing, H.264 is expected to replace H.263, as illustrated in Figure \ref{fig:Comparison-H264-H263}. It is a standard intended to cover video encoding for interactive communication, streaming and broadcasting. The main obstacle to its widespread adoption is its higher processor requirements: encoding is reported to be up to eight times more complex and decoding three times more complex than MPEG-2 (source: \cite{Latency-H264}). The transmission of H.264 over RTP is covered by RFC 3984 \cite{RFC-3984-RTP-Payload-Format-H.264}.

\begin{figure}[htp]
\centering
\includesvg[width=0.5\textwidth]{h264-vs-h263}
\caption{Comparison of the H.263 and H.264 codecs (source: \protect\cite{Comparison-H264-H263-MPEG2}).}
\label{fig:Comparison-H264-H263}
\end{figure}

\subsection{Video: MPEG-2}
This codec supports high-quality communication at bandwidths ranging from 1.5 Mbit/s to 15 Mbit/s. However, compatibility problems between different implementations prevent its widespread adoption. In terms of encoding quality, it performs worse than H.264 (Figure \ref{fig:Comparison-H264-MPEG2}). On the other hand, it requires less processing power.

\begin{figure}[htp]
\centering
\includesvg[width=0.5\textwidth]{h264-vs-mpeg2}
\caption{Comparison of the H.264 and MPEG-2 codecs (source: \protect\cite{Comparison-H264-H263-MPEG2}).}
\label{fig:Comparison-H264-MPEG2}
\end{figure}

\subsection{Video: Motion JPEG}\label{kap:MJPEG}
MJPEG encodes video frames as separate JPEG images, which are divided into RTP packets during transmission (RFC 2435 \cite{RFC-2435-RTP-Payload-Format-JPEG}). Its advantages include fast compression and decompression, as well as the ability to edit the video stream without calculating transitions between individual frames. This property is useful when processing images in real time. MJPEG also supports practically arbitrary image sizes at high quality, which is useful, for example, when transmitting desktop screenshots. Its disadvantages are higher bandwidth requirements and greater susceptibility to errors caused by poor connections and RTP packet loss.

\subsection{Audio: L16}\label{kap:L16}
A simple audio codec that supports the transmission of uncompressed audio at a resolution of 16 bits. Signal levels are represented by values from -32768 to 32767. It is used as the default codec in the RAT application (Chapter \ref{kap:RAT}). One version of this codec (16 kHz mono) is compatible with the G.711 standard \cite{ITU-G.711-PCM-of-voice}. RFC 3551 \cite{RFC-3551-RTP-Profile-AudioVideo} specifies the RTP payload profile for L16 audio at 44100 Hz with one or two channels.

\subsection{Audio: GSM}\label{kap:GSM}
GSM (Groupe Spéciale Mobile) defines the EU GSM 06.10 standard for voice encoding at a bitrate of 13,200 b/s. RFC 3551 \cite{RFC-3551-RTP-Profile-AudioVideo} specifies the RTP payload profile for GSM audio at 8000 Hz with one channel.

\subsection{Audio: G.722}\label{kap:G722}
This proprietary codec is used in hardware H.323 clients and is specified in ITU-T G.722 \cite{ITU-G.722-7KHz-Audio}.

\subsection{Audio: Speex}\label{kap:Speex}
Speex was developed as an open standard and is not encumbered by software patents. An RFC for the Speex RTP payload format has not yet been published, but an Internet-Draft \cite{IETFDraft-RTP-Payload-Speex} addresses this topic. Speex has the potential to become the most widely used codec in open-source videoconferencing applications. The performance characteristics reported in the literature \cite{Latency-AudioCodecs} support this choice.

%%%%%%%%%%%%%%%%%%%%%%%%%%%%%%%%%%%%%%%%%%%%%%%%%%%%%%%%%%%%%%%%%%%%%
%%%%%%%%%%%%%%%%%%%%%%%%%%%%%%%%%%%%%%%%%%%%%%%%%%%%%%%%%%%%%%%%%%%%%
\section{Evaluation of codec selection}
%%%%%%%%%%%%%%%%%%%%%%%%%%%%%%%%%%%%%%%%%%%%%%%%%%%%%%%%%%%%%%%%%%%%%
%%%%%%%%%%%%%%%%%%%%%%%%%%%%%%%%%%%%%%%%%%%%%%%%%%%%%%%%%%%%%%%%%%%%%

We have described several commonly used codecs suitable for videoconferencing and explained why multimedia data must be encoded during network transmission. Among video codecs, we selected H.263 as a suitable candidate because of its widespread support in existing systems, and H.264 because of its excellent compression ratio. Among audio codecs, we prefer Speex. Given the use of G.722 and H.261 in many hardware H.323 clients, support for these codecs should also be considered.

% **********************************************
% Audio Codecs with Latency up to 50ms
% **********************************************
\begin{table}
\centering
\begin{tabular}[t]{|c|c|c|c|c|c|c|c|c|}
\hline
Codec  & B/ & frm & frm/ & pak/ & pay  & B/ & bits/ & latency \\
 & frm & (ms) & pak & sec & (B) & pak & sec & \\
\hline
G723.1 & 20 & 30 & 1 & 33.33 & 20 & 60 & 16000 & 30 \\
G723.1 & 24 & 30 & 1 & 33.33 & 24 & 64 & 17067 & 30 \\
GSM & 33 & 20 & 1 & 50 & 33  & 73 & 29200 & 20 \\
GSM & 33 & 20 & 2 & 25 & 66  & 106 & 21200 & 40 \\
G.711 & 240 & 30 & 1 & 33.33 & 240 & 280 & 74667  & 30 \\
LPC10 & 7 & 22.5 & 1 & 44.44 & 7  & 47 & 16711 & 23 \\
LPC10 & 7 & 22.5 & 2 & 22.22 & 14 & 54 & 9600 & 45 \\
Speex-5.95 & 15 & 20 & 1 & 50 & 15 &  55 & 22000 &  20 \\
Speex-5.95 & 15 & 20 & 2 & 25 & 30 &  70 & 14000 &  40 \\
Speex-8k & 20 & 20 & 1 & 50 & 20  & 60 & 24000  & 20 \\
Speex-8k & 20 & 20 & 2 & 25 & 40  & 80 & 16000 &  40 \\
Speex-11k & 28 & 20 & 1 & 50 & 28 & 68 & 27200 &  20 \\
Speex-11k & 28 & 20 & 2 & 25 & 56 & 96 & 19200 &  40 \\
Speex-15k & 38 & 20 & 1 & 50 & 38 & 78 & 31200 &  20 \\
Speex-15k & 38 & 20 & 2 & 25 & 76 & 116 & 23200 &  40 \\
Speex-18.2 & 46 & 20 & 1 & 50 & 46 & 86 & 34400 &  20 \\
Speex-18.2 & 46 & 20 & 2 & 25 & 92 & 132 & 26400 &  40 \\
Speex-24.6 & 62 & 20 & 1 & 50 & 62 & 102 & 40800 &  20 \\
Speex-24.6 & 62 & 20 & 2 & 25 & 124 & 164 & 32800 &  40 \\
\hline
\end{tabular}
\caption{Audio codecs with a latency of up to 50 ms (source: \protect\cite{Latency-AudioCodecs}).}
\label{tab:AudioCodecsMinLatency}
\end{table}
% **********************************************

At the end of this chapter, we would also like to include Table \ref{tab:AudioCodecsMinLatency}, which lists audio codecs with a latency of up to 50 ms. These are low-latency codecs that could be used when collaborating on the MARCEL network project (Chapter \ref{kap:AGvUmeleckejKomunite}).
\proofed
